\section{Discusión}
El objetivo se respondió afirmativamente para estados de equilibrio condicionados: las propiedades de las especies de partícula y el exceso bariónico cosmológico permitieron inferir cantidades separadas de bariones y antibariones a una temperatura bibliográfica especificada. El resultado central a 50 MeV cuantifica una población casi simétrica y un exceso pequeño, con 306.1 antibariones por unidad neta. Al enfriarse los estados calculados, el mismo cociente alcanzó la unidad cerca del hito de 38.2 MeV publicado por \citet{Yang2022}. La sucesión coincide con el comportamiento general descrito por \citet{Satarov2013}.

La coincidencia debe interpretarse dentro de las condiciones compartidas. Ambos enfoques relacionan equilibrio hadrónico con información cosmológica sobre el exceso por entropía. Aquí se utilizaron inventarios fijados y una implementación independiente de las sumas térmicas; no se reprodujo la evolución cinética de Satarov et al. ni se realizó una réplica completa del modelo de Yang y Rafelski. Estos últimos imponen neutralidad eléctrica y de extrañeza, mientras nuestro cálculo mantuvo $\mu_Q=\mu_S=0$. Además, sus entradas cosmológicas no son idénticas. Por ello, la diferencia con 38.2 MeV no puede atribuirse exclusivamente al catálogo, ni presentarse como una prueba independiente de observaciones tempranas.

La condición química simplificada tiene un efecto verificable: sumando las cargas de las especies de partícula exportadas, a 50 MeV se obtuvo un número de carga neta igual a aproximadamente $0.4901\,\Delta n_B$ y un número de extrañeza neta de $-0.03253\,\Delta n_B$. Los leptones con potencial químico nulo no compensaron esa carga. El estado reconstruido es, por tanto, una aproximación térmica declarada; completar la reconstrucción cosmológica física requeriría resolver simultáneamente los balances de carga, extrañeza y sector leptónico. Esta limitación subsiste aunque el cruce numérico redondee a la referencia publicada.

Las cantidades comóviles aclaran la relación con el presente. El exceso se conserva bajo los supuestos adoptados, mientras el total de bariones y antibariones de equilibrio disminuye. Comparar directamente un cm$^3$ temprano y uno actual mezclaría el cambio de población con la dilución por expansión. La normalización por entropía permite comparar una misma región y declarar cantidades medias sin introducir un volumen cosmológico global no documentado. Las abundancias de H, He y D sirven como referencias nucleares externas de otra etapa; su cálculo histórico exige una red de nucleosíntesis posterior.

La reconstrucción histórica completa requiere tasas de producción inversa, aniquilación y decaimiento, junto con expansión y condiciones iniciales. En un esquema homogéneo, una ecuación del tipo
\begin{align}
\frac{dn_{\bar B}}{dt}+3Hn_{\bar B}
&=\langle\sigma v\rangle\bigl(n_B^{\rm eq}n_{\bar B}^{\rm eq}
-n_Bn_{\bar B}\bigr)\label{eq:kinetic}
\end{align}
introduce información que una tabla de equilibrio no contiene; $H$ es la tasa de expansión y $\langle\sigma v\rangle$ el promedio térmico del producto de la sección eficaz y la velocidad relativa. No se integró esa ecuación en este trabajo. Por ello no se determinaron el primer evento, el final cinético ni el número acumulado de partículas regeneradas.

La aportación del trabajo es una referencia cuantitativa reproducible que conecta entradas actuales con inventarios tempranos bajo supuestos explícitos, identifica la información recuperable y muestra el alcance del contraste publicado. No se establece una nueva observación, una explicación del origen de la asimetría o un mecanismo físico nuevo. La convergencia numérica y la sensibilidad a una entrada no sustituyen la evaluación de neutralidad, interacciones, anchos de resonancia y validez del equilibrio.
